\section{Compact spectral sets and rational descent}\label{desc:section}

We first verify the elementary properties of the sets in
\eqref{intro:compactsets}. These properties supply the uniqueness used in
Theorem~\ref{meas:detector}; their proof is independent of that theorem.

\begin{lemma}\label{desc:compactsets}
There are finitely many sets $B_o$. Each is compact and independent of
the triple and compact form used to define it, and distinct orbits give
disjoint sets. The algebra of representation characters restricted to
$B=\bigcup_oB_o$ contains the constants, separates points and is closed
under complex conjugation. It is uniformly dense in $C(B,\C)$.
Each $B_o$ is preserved by inversion and by multiplication by a
finite-order central element of $\Ghat_\C$.
\end{lemma}

\begin{proof}
Nilpotent orbits in a complex semisimple Lie algebra are finite in number.
Triples for a fixed orbit are conjugate, and maximal compact forms of a
complex reductive group, including its components, are conjugate.
Consequently their images in the affine conjugacy quotient give the same
$B_o$. The image of a compact group under the displayed continuous map is
compact. The finite union $B$, as a compact subspace of a complex affine
variety, is metrizable.

Choose a maximal torus $T\subset\Ghat_\C$. Polar decomposition in
$T(\C)$ gives
\[
T(\C)=T(\mathbb R_{>0})\times T_{\cpt},
\]
where the positive factor is intrinsically described by the absolute
values of all characters. This decomposition is Weyl-equivariant.
The centralizer of the torus $\lambda_e(\Gm)$ in $\Ghat_\C$ is a
connected reductive Levi subgroup. It contains $s$, and a maximal
torus of this Levi containing the semisimple element $s$ also contains
its central torus $\lambda_e(\Gm)$. Thus $\lambda_e$ and $s$ can be
placed in a common maximal torus. The positive part of
$\lambda_e(\sqrt r)s$ is exactly $\lambda_e(\sqrt r)$. Its Weyl orbit
determines the dominant cocharacter $\lambda_e$: evaluation of
cocharacters at $\sqrt r>1$ is injective.

For completeness, that cocharacter determines the nilpotent orbit
without forgetting any distinction in type~D. Fix $\lambda$ and write
$\ghat_i$ for its weight spaces. Whenever a triple has diagonal
cocharacter $\lambda$, the $\mathfrak{sl}_2$ decomposition gives
\[
[\ghat_0,e]=\ghat_2.
\]
Indeed, the raising map from weight zero to weight two is surjective
on each irreducible $\mathfrak{sl}_2$ summand. Thus the centralizer of
$\lambda$ has an open orbit through $e$ in the irreducible vector space
$\ghat_2$. Two such open orbits must coincide. It follows that two
parameters with the same positive part have the same ambient nilpotent
orbit. This proves disjointness; it is also the description underlying
\cite[Lemma~3.1.3]{GLR}.

Representation characters separate semisimple conjugacy values in a
connected complex reductive group: restriction to $T$ identifies their complex
span with $\C[T]^W$ \cite[Theorem~22.38]{MilneAG}.
Hence their restrictions separate points of $B$.
To check complex conjugation, write $a=\lambda_e(\sqrt r)$.
For a representation $D$, $D(a)$ is diagonalizable with positive real
eigenvalues and commutes with $D(s)$; on each of its eigenspaces $D(s)$
is unitary after choosing an invariant Hermitian form for the compact
centralizer. Therefore
\[
\overline{\chi_D(as)}=\chi_D(as^{-1}).
\]
A Weyl element in the triple's $\mathrm{SL}_2$ conjugates $a$ to
$a^{-1}$ and commutes with $s$. Thus
\[
\chi_D(as^{-1})=\chi_D(a^{-1}s^{-1})=\chi_{D^\vee}(as).
\]
The same representation $D^\vee$ works simultaneously on every $B_o$.
The character algebra is consequently self-adjoint, and the complex
Stone--Weierstrass Theorem proves its density on the finite compact union.

The triple Weyl element also shows that $(as)^{-1}$ is conjugate to
$as^{-1}$, proving invariance under inversion. A finite-order central
element lies in the chosen compact triple centralizer: adjoining it
gives a compact group, and maximality leaves that group unchanged.
Multiplication by it therefore preserves $B_o$.
\end{proof}

\begin{corollary}\label{desc:unique-measure}
A finite complex Borel measure on $B$ is determined by its integrals
against all representation characters.
\end{corollary}

\begin{proof}
Integration against a finite complex measure is a continuous functional
on $C(B,\C)$ for the uniform norm. Lemma~\ref{desc:compactsets} makes the
character algebra dense. Equality on the characters therefore gives
equality of the continuous functionals and hence of the measures.
\end{proof}

\subsection{Independence of the coefficient isomorphism}

All algebraically closed characteristic-zero coefficient fields below
are considered as abstract fields. Both $E$ and $\C$ have transcendence
degree $2^{\aleph_0}$ over $\Q$, so there are field isomorphisms
$E\simeq\C$. No continuity or compatibility with a selected embedding of
the algebraic numbers is used.

\begin{proposition}\label{desc:coefficient-independence}
For a fixed closed invariant support $Y$, the subspace
$\iota_*(\cF_Y\cA_E)\subset\cA_\C$ is independent of the abstract
isomorphism $\iota:E\xrightarrow{\sim}\C$, after matching nilpotent
orbits by the split root datum. Consequently $\cF_Y\cA_E$ is invariant
under every automorphism of $E$ fixing $\Q$.
\end{proposition}

\begin{proof}
The point functions identify every transported ambient space with the
same $\C^{(S)}$. The pinned split form identifies representation
characters by highest weight across coefficient fields. The pairing \eqref{intro:pairing} has rational
coefficients. Spherical Hecke correspondences count the same finite-field
modifications under each transport. Their normalization may be chosen
on the complex side using the fixed positive $\sqrt q$, so their
operators $T_D$ are the same for each $\iota$.

If a transported choice of $\sqrt q$ differs from this one, set
\[
z=(2\rho_G)(-1)\in Z(\Ghat)[2],
\]
where $2\rho_G$ is the sum of the positive roots of $G$, viewed as a
cocharacter of its dual group. The Satake parity convention at the
degree-$d$ point introduces $z^d$: a highest weight $\mu$ acquires the
sign $(-1)^{d\langle 2\rho_G,\mu\rangle}$, which is exactly its central
character at $z^d$. This convention is also encoded by the modified
dual group in \cite[Remark~6.7 and Section~6.4]{RicharzScholbach}.
On characters the change is therefore translation by $z^d$, and on
moment measures it is the corresponding pushforward. By
Lemma~\ref{desc:compactsets}, this translation preserves each $B_o$.
The contragredient convention similarly preserves each $B_o$ by
inversion. Thus these conventions do not change the support test.

Nilpotent orbits of the split group are identified across these fields
using the fixed pinned $\Q$-form and their triple cocharacters, or
equivalently their weighted Dynkin diagrams. No outer automorphism is
introduced in this identification. In particular the orbit labels in
$Y$ and the compact set
$B_Y$ are fixed independently of $\iota$. Theorem~\ref{meas:detector}
and Corollary~\ref{desc:unique-measure} now characterize the transported
subspace by the same moment equations and the same support condition
for every test $g\in\cA_\C$. The subspaces are equal.

Fix one $\iota$ and let $\tau\in\Aut(E/\Q)$. Comparing $\iota$ with
$\iota\circ\tau$ gives
$\iota_*(\cF_Y\cA_E)=\iota_*\tau_*(\cF_Y\cA_E)$.
Applying $\iota^{-1}$ proves invariance.
\end{proof}

\subsection{Descent on finite coordinate sets}

\begin{lemma}\label{desc:descent}
Let $E$ be an algebraically closed extension of $\Q$, and let
$V_\Q$ have a specified basis, possibly infinite. If
$W\subset E\otimes_\Q V_\Q$ is invariant under every
$\tau\in\Aut(E/\Q)$, then
\[
E\otimes_\Q(W\cap V_\Q)\longrightarrow W
\]
is an isomorphism.
\end{lemma}

\begin{proof}
The fixed field of $\Aut(E/\Q)$ is $\Q$. Indeed, an algebraic element
outside $\Q$ has a different conjugate, and the corresponding embedding
of its number field extends to an automorphism of $E$. A transcendental
element can be included in a transcendence basis; the substitution
$t\mapsto t+1$ on that basis element extends from the rational function
field to its algebraic closure $E$. Thus it too is moved by an
automorphism over $\Q$.

Let $I$ be a finite subset of the specified basis, choose an order on
$I$, and put
\[
W_I=W\cap E^I.
\]
This is an invariant finite-dimensional subspace. Its unique reduced
row-echelon basis matrix is unchanged by every automorphism over $\Q$:
field automorphisms preserve the row space, the pivot positions and the
normalization conditions. Every matrix entry therefore belongs to the
fixed field $\Q$. Its nonzero rows give a rational basis, and hence
\[
W_I=E\otimes_\Q(W_I\cap\Q^I).
\]
Every vector of $E\otimes_\Q V_\Q$ has finite support in the specified
basis. Taking the directed union over $I$ gives the claimed
surjectivity. Injectivity follows by tensoring the inclusion
$W\cap V_\Q\hookrightarrow V_\Q$ with the flat $\Q$-module $E$.
\end{proof}

\begin{proof}[Proof of Theorem~\ref{intro:main}]
Proposition~\ref{desc:coefficient-independence} makes
$W=\cF_Y\cA_E$ invariant under all coefficient automorphisms over $\Q$.
Apply Lemma~\ref{desc:descent} to $V_\Q=\cA_\Q$ with its point-function
basis. This is exactly the map \eqref{intro:main-map}.

The argument applies to every closed invariant $Y$, including the empty
set and the full nilpotent cone. The trace calculation uses all retained
fixed components and actual finite linear combinations of Weil trace
classes. The detector allows every finitely supported test function.
Consequently the result holds for the entire specified space, without
any additional finiteness or cuspidality assumption.
\end{proof}

The argument also compares coefficient primes on the full function space
under the same four conditions of Assumption~\ref{intro:characteristic}.
Put $p=\operatorname{char}(\mathbb F_q)$ and
$E_\ell=\overline{\Q}_\ell$ for $\ell\ne p$. Let
$Y_\ell\subset\Nhat_{E_\ell}$ be closed invariant supports with the same
nilpotent-orbit labels under the pinned split dual group. Then, as
subspaces of $\cA_\Q$,
\[
 \cA_\Q\cap\cF_{Y_\ell}\cA_{E_\ell}
 =\cA_\Q\cap\cF_{Y_{\ell'}}\cA_{E_{\ell'}}
 \qquad(\ell,\ell'\ne p).
\]
Indeed, choose the same closed point $v$ and abstract isomorphisms
$E_\ell\simeq\C$. With the positive square-root Satake normalization,
the moment equations for every $f,g\in\cA_\C$ and the compact support
sets in Theorem~\ref{meas:detector} are the same for every prime, as in
the proof of Proposition~\ref{desc:coefficient-independence}. Uniqueness
of the measures gives the same transported filtration in $\cA_\C$.
Intersecting with $\cA_\Q$, which every coefficient isomorphism fixes,
proves the equality. As the matched closed orbit subset varies, these
common rational subspaces form a filtration of $\cA_\Q$ whose
$E_\ell$ scalar extension is the original Arthur filtration, by
Theorem~\ref{intro:main}. This compares rational function subspaces;
it does not identify spectral support categories or individual
eigenobjects.

\subsection{Unramified cuspidal comparison}

\begin{corollary}[Cuspidal comparison across coefficient primes]
\label{desc:cuspidal-comparison}
Keep the curve $X/\mathbb F_q$, the split connected semisimple group
$G/\mathbb F_q$, and all four conditions of
Assumption~\ref{intro:characteristic}. Put
$p=\operatorname{char}(\mathbb F_q)$ and
$E_\ell=\overline{\Q}_\ell$ for $\ell\ne p$. At unramified level $D=0$,
let $C_{0,\Q}\subset\cA_\Q$ be the cuspidal function space and
$C_{0,\mathcal O,\Q}$ its rational orbit summands from
\cite[Theorem~1.1]{CuspidalArthur}, and set
$C_{0,\ell}=E_\ell\otimes_\Q C_{0,\Q}$.
For each $\ell\ne p$, let $Y_\ell\subset\Nhat_{E_\ell}$ be closed and
invariant, with the same nilpotent-orbit subset $\mathscr Y$ under the
fixed pinned split rational dual group. Write $Y=(Y_\ell)_{\ell\ne p}$ and
define
\[
 F^{\Q}_{Y,\ell}:=\cA_\Q\cap\cF_{Y_\ell}\cA_{E_\ell},
 \qquad
 V_{Y,\Q}:=\bigoplus_{\mathcal O\in\mathscr Y}C_{0,\mathcal O,\Q}.
\]
Then
\begin{equation}\label{desc:cuspidal-equalities}
 \begin{split}
 C_{0,\Q}\cap F^{\Q}_{Y,\ell}&=V_{Y,\Q},\\
 C_{0,\ell}\cap\cF_{Y_\ell}\cA_{E_\ell}
   &=E_\ell\otimes_\Q V_{Y,\Q}.
 \end{split}
\end{equation}
The intersections are taken in the corresponding finitely supported
function spaces. The equalities include the zero cusp space and empty or
full $Y_\ell$.
\end{corollary}

This is precisely the unramified cuspidal common-scope form of
\cite[Conjectures~3.4.9 and~3.4.11]{GLR}. The additional identification
with the companion's orbit summands concerns the $D=0$ cuspidal subspace.

\begin{proof}
For split $G$, the companion's unramified double quotient is the set $S$
of bundle classes \cite[Section~6.1]{CuspidalArthur}, giving the stated
point-function embeddings. Fix $\ell$, an abstract isomorphism
$\iota:E_\ell\xrightarrow{\sim}\C$, and a closed point $v$ with
$r=q^{\deg v}$. The finite reduced global Hecke image on the cusp space
gives finitely many joint eigenspaces \cite[Lemma~6.1]{CuspidalArthur}.
Restrict to the spherical algebra at $v$ and group all joint characters
with the same one-place Satake class. The ordinary Hecke compatibility of
Theorem~\ref{inp:trace} and the companion's volume-one normalization give
\[
 C_{0,\C}:=\C\otimes_\Q C_{0,\Q}
   =\bigoplus_{t\in\Sigma_v}C_t,
 \qquad T_V|_{C_t}=\chi_V(t),
\]
where the $t$ are distinct classes in the positive $\sqrt r$ convention
and $V$ is any finite-dimensional representation of $\Ghat_\C$.
The possible square-root correction is a finite central translate
\cite[Remark~6.3]{CuspidalArthur}, and a contragredient convention gives
inversion; both preserve every $B_{\mathcal O}$ by
Lemma~\ref{desc:compactsets}.

The companion's orbit descent assigns each joint character one orbit
label, independent of the good place and complex embedding, and the complex
scalar extension of each rational orbit summand is the sum of the joint
eigenspaces with that label
\cite[Lemma~6.2 and Propositions~6.5--6.6]{CuspidalArthur}. Every occurring
embedded class of label $\mathcal O$ lies in $B_{\mathcal O}$ by the
all-embeddings compactness defining $R_{r,\mathcal O}$. Since the
$B_{\mathcal O}$ are disjoint, all joint eigenspaces grouped into one
$C_t$ have the same label. Thus $\Sigma_v\subset B$ and
\[
 \C\otimes_\Q C_{0,\mathcal O,\Q}
   =\bigoplus_{t\in\Sigma_v\cap B_{\mathcal O}}C_t.
\]
For the common support put
$B_Y=\bigcup_{\mathcal O\in\mathscr Y}B_{\mathcal O}$.

Write $f=\sum_t f_t\in C_{0,\C}$. For every $g\in\cA_\C$, the finite
atomic measure $\mu_{f,g}:=\sum_t[f_t,g]\delta_t$ on $B$ satisfies
\[
 [T_Vf,g]=\sum_t\chi_V(t)[f_t,g]
           =\int_B\chi_V\,d\mu_{f,g}.
\]
Corollary~\ref{desc:unique-measure} identifies $\mu_{f,g}$ with the
detector measure $\nu_{f,g}$. If $f_t\ne0$, choose $b\in S$ with
$f_t(b)\ne0$ and take the point function $g=\delta_b$. Then
\[
 \mu_{f,\delta_b}(\{t\})=[f_t,\delta_b]
             =\frac{f_t(b)}{|\Aut(b)|}\ne0.
\]
The atoms are distinct, so no other local component cancels this mass.
The test need not be cuspidal. Consequently all detector measures are
supported in $B_Y$ exactly when $f_t=0$ for every $t\notin B_Y$.
Theorem~\ref{meas:detector} and the preceding orbit decomposition give
\[
 C_{0,\C}\cap\iota_*(\cF_{Y_\ell}\cA_{E_\ell})
   =\C\otimes_\Q V_{Y,\Q}.
\]
Pulling back by $\iota$, which fixes $\Q$, proves the second equality in
\eqref{desc:cuspidal-equalities}. Intersect it with $C_{0,\Q}$ and use
$C_{0,\Q}\cap(E_\ell\otimes_\Q V_{Y,\Q})=V_{Y,\Q}$ to obtain the first.
\end{proof}
